\documentclass[12pt]{article}

\usepackage{amsmath}

\setlength{\textwidth}{6.5in}
\setlength{\textheight}{7.5in}
\setlength{\hoffset}{-0.75in}
\setlength{\voffset}{-0.75in}

\setlength{\parindent}{0.0in}
\setlength{\parskip}{0.5\baselineskip}

\input{latex_macros.tex}

\title{Constructing Diversification Constraints}
\author{R. Scott McIntire}
\date{Aug 11, 2026}

\begin{document}

\maketitle

\section{Overview}
When one constructs (or modifies) a long-only portfolio of $n$ securities
--- a vector ${\bf x} \succeq {\bf 0}$ of notional positions%
\footnote{A bold font is used to denote vectors. We refer to a
vector of $0$s as ${\bf 0}$, and likewise we refer to a vector of $1$s
as ${\bf 1}$. The mathematical expression, ${\bf a} \preceq {\bf b}$, is to be
interpreted as the statement that every element of the vector ${\bf a}$ is
less than or equal to the corresponding element in the vector ${\bf b}$.
Similarly, the expression, ${\bf a} \succeq {\bf b}$, is the statement that
each element of ${\bf a}$ is greater than or equal to each
corresponding element of ${\bf b}$.
}
--- a diversification mandate will often demand that the largest positions
not dominate the book: {\it the $k$ largest positions may not, together,
exceed a fraction $p$ of the portfolio's total notional}.%
\footnote{Rules of this shape occur in practice; the European UCITS
``5/10/40'' rule is a well-known relative.}
The difficulty in enforcing such a rule inside an optimization is that
{\it which\/} $k$ positions are the largest depends on ${\bf x}$ --- the
very variable being optimized. One way out is to write down, for every
set of $k$ securities, a constraint bounding the sum of those positions;
but the number of constraints that one has to write is $\binom{n}{k}$.
This becomes combinatorially large very quickly. In this paper we use LP
duality to replace these $\binom{n}{k}$ constraints with $O(n)$ linear
constraints.

To formalize, let $k \in {\bf Z}^+$ with $k \le n$ throughout, and
consider the function, $f: {\bf R}^n \to {\bf R}$, defined by%
\footnote{
For a given vector, ${\bf x} \in {\bf R}^n$, create a vector ${\bf x}^\prime$
by sorting ${\bf x}$ from highest to lowest (ties broken arbitrarily, as the
sums below are independent of the tie-breaking rule).
Define the notation, $x_{[i]}$ by: $x_{[i]} \equiv x_{i}^\prime$.}
\begin{align*}
	f({\bf x}) &= \sum_{i=1}^k x_{[i]}.
\end{align*}
That is, $f({\bf x})$ is the sum of the $k$ largest entries of ${\bf x}$.
Note that $f$ is {\it convex\/}: it is the maximum, over all $\binom{n}{k}$
index subsets $S$ of size $k$, of the linear functions
${\bf x} \mapsto \sum_{i \in S} x_i$.  This convexity is the reason the
constraint $f({\bf x}) \le M$ can admit a linear representation at all.
We are interested in optimization problems involving a vector ${\bf x}$
where the sum of the top $k$ entries is bounded; that is, $f({\bf x}) \le M$
for a given $M$.  The portfolio mandate is the case
$M = p \, {\bf 1}^T {\bf x}$.  This bound depends on the variable
${\bf x}$ --- but only {\it linearly\/}, so the constraint (\ref{opp1})
below remains linear with $M$ replaced by $p \, {\bf 1}^T {\bf x}$.
The function $f$ is also $k$ times the
{\it conditional value-at-risk\/} (CVaR) of the empirical distribution of
the entries of ${\bf x}$ at level $k/n$; the reformulation below is
precisely the linearization of Rockafellar and
Uryasev~\cite{rockafellar-uryasev} in that setting.

This reformulation is classical --- it appears in Ogryczak and
Tamir~\cite{ogryczak-tamir}; see also Boyd and
Vandenberghe~\cite{boyd} --- but the derivation below is self-contained.
The route is a {\it change of representation\/}, made three times: the
sorted sum becomes a $0$--$1$ integer program; the integer program
becomes a continuous linear program; and the linear program is traded for
its {\it dual\/}.  Each trade is forced by a defect of {\it shape\/} in
the current description; the table at the end of the paper collects the
ladder.

\section{The Function $f$ as an Optimization Problem}
We wish to characterize $f({\bf x})$ as the value of an optimization.
We claim that for a given ${\bf x}$, $f({\bf x})$ equals the optimal value of
the following constrained {\it discrete\/} optimization problem:%
\begin{align}
	\max_{{\bf y} \in \{0,1\}^n} \quad & {\bf y}^T {\bf x} \label{discrete_objective} \\
	\text{s.t.} \quad & {\bf y}^T {\bf 1} = k \label{discrete_card}
\end{align}
This is $f({\bf x})$ written as a 0-1 integer program (IP).
Although succinct, this optimization still requires examining $\binom{n}{k}$
candidate solutions.

\textbf{Claim.} The discrete optimization problem is equivalent to the following continuous problem:%
\footnote{%
The objective ${\bf y}^T {\bf x}$ is continuous and the feasible set is closed and bounded
(a compact polytope in ${\bf R}^n$); the supremum is therefore attained, so we may
write ``max'' instead of ``sup''.
}
\begin{align}
	\max_{{\bf y} \in {\bf R}^n} \quad & {\bf y}^T {\bf x} \notag \\
	\text{s.t.} \quad & {\bf 0} \preceq {\bf y} \preceq {\bf 1} \label{bound} \\
	& {\bf y}^T {\bf 1} = k \label{cont_card}
\end{align}
\textbf{Proof.}
Since its feasible set contains the discrete problem's feasible set
(every ${\bf y} \in \{0,1\}^n$ with ${\bf y}^T {\bf 1} = k$ also satisfies
${\bf 0} \preceq {\bf y} \preceq {\bf 1}$), the optimal value of the continuous problem is
at least as large as the optimal value of the original discrete problem.
We now show the reverse. For any optimal solution ${\bf y}^*$ of the continuous
problem, we construct a ${\bf y}^\prime \in \{0,1\}^n$ satisfying
(\ref{discrete_card}) and yielding the same value of the objective
(\ref{discrete_objective}) as ${\bf y}^*$.

If all components of ${\bf y}^*$ are integers, by (\ref{bound}) they are each
$0$ or $1$; in this case take ${\bf y}^\prime = {\bf y}^*$.

Otherwise, consider the components of ${\bf y}^*$
that are non-integer; that is, have values strictly between $0$ and $1$.
By (\ref{cont_card}), $\sum_i y^*_i = k$ is an integer; the integer-valued
components contribute an integer to this sum, so the non-integer components
must sum to an integer too. Consequently, there must be more than one such
component. It must also be the case that amongst these non-integer components
the corresponding entries of ${\bf x}$ are the same. Otherwise, there would
be indices $i, j$ with $y^*_i, y^*_j \in (0,1)$ and $x_i > x_j$; shifting a
small amount $\epsilon > 0$ from $y^*_j$ to $y^*_i$ preserves the
constraints (\ref{bound}) and (\ref{cont_card}) (both components are
strictly between $0$ and $1$, so a small shift stays within the bounds) and
increases the objective by $\epsilon (x_i - x_j) > 0$ --- contradicting the
optimality of ${\bf y}^*$. Therefore, all of the entries of ${\bf x}$
associated with non-integer components of ${\bf y}^*$ have the same
value. As previously stated, the sum of these non-integer values is an integer,
which is positive and necessarily less than the number of non-integer
components in ${\bf y}^*$ (each such component lies strictly between $0$ and $1$).
Let $j$ denote this integer sum.
Pick any $j$ of these non-integer values, and let $J$ be the set of their index
values in ${\bf y}^*$.
Let ${\tilde J}$ be the set of indices corresponding to the remaining non-integer values
of ${\bf y}^*$.

Now, create a new 
vector ${\bf y}^\prime \in \{0,1\}^n$ and set its values in the following way: 
Set the values of ${\bf y}^\prime$ corresponding to the indices in $J$
to be $1$. Set the values of ${\bf y}^\prime$ corresponding to the indices in ${\tilde J}$
to $0$. Set the rest of the components of ${\bf y}^\prime$ to the
corresponding values from ${\bf y}^*$. By our construction,
${{\bf y}^\prime}^T {\bf 1} = k$. The components of ${\bf y}^\prime$ inherited from ${\bf y}^*$
were already integers, and by (\ref{bound}) every integer component of ${\bf y}^*$
is $0$ or $1$. It follows that ${\bf y}^\prime \in \{0,1\}^n$ and
satisfies the constraint of the discrete problem, (\ref{discrete_card}).
The objective value is also preserved: let $x_c$ denote the common value of
the entries of ${\bf x}$ on $J \cup {\tilde J}$. The contribution to the objective from those indices
equals $x_c \sum_{i \in J \cup {\tilde J}} y^*_i = x_c \, j$, which is exactly the
contribution from ${\bf y}^\prime$ on the same indices.
Therefore, ${\bf y}^\prime$ is in the feasible set of the discrete problem and
attains the optimal value of the continuous problem; that is, the
optimal value of the discrete problem is at least the optimal value of the
continuous problem. Consequently, the two optimization problems are equivalent.

\textbf{Remark.} The equivalence is no accident: the vertices of the polytope
$\{ {\bf y} : {\bf 0} \preceq {\bf y} \preceq {\bf 1},\; {\bf y}^T {\bf 1} = k \}$
are exactly the $0$--$1$ vectors with $k$ ones, and a linear objective always
attains its maximum at a vertex.  The proof above has the advantage of being
self-contained.

The continuous problem above is equivalent to the following minimization:
\begin{align}
	\min_{{\bf y} \in {\bf R}^n} \quad & -{\bf y}^T {\bf x} \label{op1} \\
	\text{s.t.} \quad & {\bf 0} \preceq {\bf y} \preceq {\bf 1} \label{op2} \\
	& {\bf y}^T {\bf 1} = k \label{op3}
\end{align}
Treating ${\bf x}$ as a constant, this is a {\it convex\/} optimization problem.

\section{A Dual Description of the Optimization}
Pause to ask what shape of description we {\it want\/}.  Every
description of $f({\bf x})$ so far is a {\it maximum\/}.  A maximum is
the wrong shape to sit inside the constraint $f({\bf x}) \le M$: to
certify that a maximum is at most $M$ one must check {\it every\/}
candidate --- for us, all $\binom{n}{k}$ of them (the last section makes
this precise).  What we want is a description of $f({\bf x})$ as a
{\it minimum\/}: to certify that a minimum is at most $M$ it suffices to
exhibit {\it one\/} feasible point with objective value at most $M$.
Duality is precisely the machine that converts the one into the other.

Since the problem described by (\ref{op1})--(\ref{op3}) is
{\it convex\/}, the value of its solution is the same as the value of
its associated {\it dual\/} problem.%
\footnote{Normally one needs to show that a convex problem satisfies Slater's condition
	in order to invoke strong duality.
For linear programs, strong duality holds whenever the primal is feasible and bounded
-- no Slater check needed.
The bilinear term $-{\bf y}^T {\bf x}$ becomes linear in ${\bf y}$ once we fix
${\bf x}$ as a parameter.}

To form the dual problem we need the Lagrangian,%
\footnote{In companion papers --- ``Properties of a Linear Unbiased
Minimum Variance Estimator'', and the minimum-norm lemma of ``Singular
Transformations of Probability Density Functions'' --- Lagrange
multipliers are scaffolding: introduced to locate a minimizer, then
discarded.  Here the multipliers are the {\it deliverable\/}:
${\boldsymbol \lambda}$ and $\nu$ survive into the final answer as the
new decision variables of the next section, and $\nu$ will even acquire
a meaning --- see the Remark below.}
which is:%
\footnote{${\bf x}$ is a parameter (constant for the optimization) -- separated from
the optimization variables by a semi-colon.}
\begin{align*}
	L({\bf y}, {\boldsymbol \lambda}_1, {\boldsymbol \lambda}_2, \nu; {\bf x}) &=
	-{\bf y}^T {\bf x} - {\boldsymbol \lambda}_1^T {\bf y}
	+ {\boldsymbol \lambda}_2^T ({\bf y} - {\bf 1}) + \nu ({\bf y}^T {\bf 1} - k).
\end{align*}
Here ${\boldsymbol \lambda}_1, {\boldsymbol \lambda}_2 \in {\bf R}^n$ with
${\boldsymbol \lambda}_1, {\boldsymbol \lambda}_2 \succeq {\bf 0}$ (sign constraints
from the two inequality systems) and $\nu \in {\bf R}$ free (from the equality)
are the {\it dual\/} variables.

Define $g$ by:
\begin{align*}
	g({\boldsymbol \lambda}_1, {\boldsymbol \lambda}_2, \nu; {\bf x}) &=
		\inf_{\bf y} \; L({\bf y}, {\boldsymbol \lambda}_1, {\boldsymbol \lambda}_2, \nu; {\bf x}).
\end{align*}

Substituting for $L$ this becomes
\begin{align*}
	g({\boldsymbol \lambda}_1, {\boldsymbol \lambda}_2, \nu; {\bf x}) &=
		\inf_{\bf y} \Bigl[\, {\bf y}^T \bigl( -{\bf x} - {\boldsymbol \lambda}_1 + {\boldsymbol \lambda}_2 + \nu {\bf 1} \bigr)
		- {\bf 1}^T {\boldsymbol \lambda}_2 - \nu \, k \,\Bigr].
\end{align*}

The dual problem is then
\begin{align*}
	\max_{\substack{{\boldsymbol \lambda}_1 \succeq {\bf 0} \\ {\boldsymbol \lambda}_2 \succeq {\bf 0} \\ \nu}} \quad
		g({\boldsymbol \lambda}_1, {\boldsymbol \lambda}_2, \nu; {\bf x}).
\end{align*}

Which is%
\footnote{Note that $g$ is $-\infty$ unless the coefficient of ${\bf y}$ is the zero vector.
Consequently, the maximum must necessarily occur where that coefficient is zero.}
\begin{align*}
	\max_{\substack{{\boldsymbol \lambda}_1 \succeq {\bf 0} \\ {\boldsymbol \lambda}_2 \succeq {\bf 0} \\ \nu}} \quad & -{\bf 1}^T {\boldsymbol \lambda}_2 - \nu \, k \\
	\text{s.t.} \quad & -{\boldsymbol \lambda}_1 + {\boldsymbol \lambda}_2 + \nu {\bf 1} = {\bf x}
\end{align*}

This is equivalent to%
\footnote{We can remove $n$ variables and the $n$ equality constraints by eliminating
	${\boldsymbol \lambda}_1$. The equality $-{\boldsymbol \lambda}_1 + {\boldsymbol \lambda}_2 + \nu {\bf 1} = {\bf x}$
	together with ${\boldsymbol \lambda}_1 \succeq {\bf 0}$ is equivalent to
	${\boldsymbol \lambda}_2 + \nu {\bf 1} \succeq {\bf x}$.
	Since there is now only one ${\boldsymbol \lambda}$, we relabel ${\boldsymbol \lambda}_2$ as ${\boldsymbol \lambda}$.}
\begin{align*}
	\max_{\substack{{\boldsymbol \lambda} \succeq {\bf 0} \\ \nu}} \quad & -\nu \, k - {\bf 1}^T {\boldsymbol \lambda} \\
	\text{s.t.} \quad & {\bf x} \preceq {\boldsymbol \lambda} + \nu {\bf 1}
\end{align*}
Negating the objective converts the max to a min:
\begin{align*}
	\min_{{\boldsymbol \lambda}, \nu} \quad & \nu \, k + {\bf 1}^T {\boldsymbol \lambda} \\
	\text{s.t.} \quad & {\bf x} \preceq {\boldsymbol \lambda} + \nu {\bf 1} \\
	& {\boldsymbol \lambda} \succeq {\bf 0}
\end{align*}

\textbf{Remark.} The dual has a useful closed form.  For fixed $\nu$, the
best choice of ${\boldsymbol \lambda}$ is clearly
$\lambda_i = (x_i - \nu)_+ \equiv \max(x_i - \nu,\, 0)$, so the dual value is
\begin{align*}
	f({\bf x}) &= \min_{\nu \in {\bf R}} \;
	\Bigl[\, k \nu + \sum_{i=1}^n (x_i - \nu)_+ \Bigr],
\end{align*}
and the minimum is attained at $\nu^* = x_{[k]}$, the $k^{\rm th}$ largest
entry of ${\bf x}$.  Indeed, at $\nu = x_{[k]}$ the bracket evaluates to
$k \, x_{[k]} + \sum_{i=1}^{k} \bigl(x_{[i]} - x_{[k]}\bigr)
= \sum_{i=1}^k x_{[i]} = f({\bf x})$, which exhibits the optimal dual point
directly, without appeal to LP theory.  The multiplier $\nu$ acts as a
threshold, and ${\boldsymbol \lambda}$ pays for the excess of each entry
above it.

The Remark buys more than an interpretation.  Weak duality is here a
one-line computation: if ${\bf x} \preceq {\boldsymbol \lambda} + \nu {\bf 1}$
with ${\boldsymbol \lambda} \succeq {\bf 0}$, and $S$ is an index set of
the $k$ largest entries of ${\bf x}$, then
\begin{align*}
	f({\bf x}) = \sum_{i \in S} x_i
	\;\le\; \sum_{i \in S} (\lambda_i + \nu)
	\;=\; \nu \, k + \sum_{i \in S} \lambda_i
	\;\le\; \nu \, k + {\bf 1}^T {\boldsymbol \lambda}.
\end{align*}
So every feasible $({\boldsymbol \lambda}, \nu)$ bounds $f({\bf x})$ from
above, and the point exhibited above attains the bound.  Together the two
lines prove --- elementarily --- the strong duality that the next section
uses.

\textbf{A calibration.} Take $n = 4$, ${\bf x} = (3, 1, 2, 2)^T$, and
$k = 2$.  Sorting gives $f({\bf x}) = 3 + 2 = 5$.  The recipe above says
$\nu^* = x_{[2]} = 2$ and
${\boldsymbol \lambda}^* = \bigl((3-2)_+,\, (1-2)_+,\, (2-2)_+,\, (2-2)_+\bigr)^T
= (1, 0, 0, 0)^T$,
with dual value $\nu^* k + {\bf 1}^T {\boldsymbol \lambda}^* = 4 + 1 = 5$.
The new machine reproduces the answer the reader already owns by sorting
--- which is how one learns to trust a new tool.  The example also
realizes the tie that the equivalence proof of Section 2 labors over:
since $x_3 = x_4 = 2$, the continuous problem has optimal solutions
${\bf y}^* = (1, 0, t, 1-t)^T$ for every $t \in [0, 1]$, all of value
$5$; the proof's rounding sends each of them to $(1,0,1,0)^T$ or
$(1,0,0,1)^T$.

\section{Characterization of $f$ as $O(n)$ Linear Constraints}
We now apply the results of the previous sections to the following problem:
Add constraints to an optimization problem that includes a {\it variable} ${\bf x} \in {\bf R}^n$
so that $f({\bf x}) \le M$ for a given $M$.

We claim that one can bound $f({\bf x})$ by adding $n+1$ new variables
(a vector ${\boldsymbol \lambda} \in {\bf R}^n$ and a scalar $\nu \in {\bf R}$),
along with the following:
\begin{align}
	\nu \, k + {\bf 1}^T {\boldsymbol \lambda} &\le M \label{opp1} \\
	{\bf x} &\preceq {\boldsymbol \lambda} + \nu {\bf 1} \label{opp2} \\
	{\boldsymbol \lambda} &\succeq {\bf 0} \label{opp3}
\end{align}

Why? Because for any $({\boldsymbol \lambda}, \nu)$ satisfying (\ref{opp2}) and (\ref{opp3}),
$\nu \, k + {\bf 1}^T {\boldsymbol \lambda}$ is an upper bound on $f({\bf x})$ -- this is weak duality.
Consequently, we have the inequality: 
$f({\bf x}) \le \nu \, k + {\bf 1}^T {\boldsymbol \lambda} \le M$.
The only concern remaining is that there may be a gap between the values of 
$\nu \, k + {\bf 1}^T {\boldsymbol \lambda}$ 
and $f({\bf x})$ -- making (\ref{opp1}) too restrictive. 
But this is not the case: by strong duality (proved elementarily in the
Remark of the previous section), the minimum of
$\nu \, k + {\bf 1}^T {\boldsymbol \lambda}$ over ${\boldsymbol \lambda}, \nu$
(subject to (\ref{opp2}) and (\ref{opp3})) equals $f({\bf x})$.
So there always exist ${\boldsymbol \lambda}, \nu$ feasible for (\ref{opp2}) and (\ref{opp3})
with $\nu \, k + {\bf 1}^T {\boldsymbol \lambda} = f({\bf x})$, and (\ref{opp1}) is satisfiable
for ${\bf x}$ exactly when $f({\bf x}) \le M$ -- the new constraints introduce no
restrictions on ${\bf x}$ beyond $f({\bf x}) \le M$ itself.

One might ask why we couldn't do exactly the same procedure with the ``simpler''
characterization we had with equations
(\ref{op1})--(\ref{op3}), rather than going through so much
effort with a dual problem. It would require the same number of variables and constraints.

We could try that; however, note that once we take the characterization of a
bound on a ``given'' ${\bf x}$ and place it into an optimization problem, that
``given'' ${\bf x}$ is no longer a constant. It is now part of the
optimization. The primal characterizes $f({\bf x})$ as a {\it maximum\/} over ${\bf y}$;
to enforce $f({\bf x}) \le M$ we would need ${\bf y}^T {\bf x} \le M$ for {\it every\/}
feasible ${\bf y}$ -- which is exactly the original combinatorial blow-up. The only way
to collapse it into a finite set of constraints is to encode the LP's optimality
conditions, which introduce bilinear terms in the (now) {\it variable\/} ${\bf x}$
and the new {\it variable\/} ${\bf y}$. These terms are decidedly {\it non-linear\/};
in fact, they are not even {\it convex\/}.

In contrast to this ``simpler'' characterization and the original na\"\i ve
combinatorial approach, we have shown that
by examining an associated dual problem, $f({\bf x})$ can be bounded by
adding $(n+1)$ new variables and $(2n + 1)$ {\it linear\/} constraints.

Return, finally, to the portfolio mandate that opened this paper.  To
require that the top $k$ positions of a long-only portfolio not exceed
the fraction $p$ of its total notional, adjoin the variables
${\boldsymbol \lambda} \in {\bf R}^n$ and $\nu \in {\bf R}$ to the
portfolio problem and impose (\ref{opp1})--(\ref{opp3}) with
$M = p \, {\bf 1}^T {\bf x}$:
\begin{align*}
	\nu \, k + {\bf 1}^T {\boldsymbol \lambda} \le p \, {\bf 1}^T {\bf x},
	\qquad
	{\bf x} \preceq {\boldsymbol \lambda} + \nu {\bf 1},
	\qquad
	{\boldsymbol \lambda} \succeq {\bf 0}.
\end{align*}
Every constraint is linear in $({\bf x}, {\boldsymbol \lambda}, \nu)$, so
any LP or QP solver accepts them unchanged, and the mandate is enforced
exactly --- no sorting, and no enumeration of $\binom{n}{k}$ subsets.

In summary, the paper has been one number wearing five descriptions, each
traded for the next because of a defect of shape:
\begin{center}
\begin{tabular}{l|l|l}
description of $f({\bf x})$ & shape & defect forcing the next rung \\ \hline
sum of the $k$ largest entries & a sort & a sort is not a constraint \\
$\max_{|S|=k} \sum_{i \in S} x_i$ & max of $\binom{n}{k}$ linear pieces & too many pieces \\
$0$--$1$ program (\ref{discrete_objective})--(\ref{discrete_card}) & max over a discrete set & combinatorial search \\
linear program (\ref{bound})--(\ref{cont_card}) & max over a polytope & a max needs {\it every\/} ${\bf y}$ checked \\
dual program (\ref{opp1})--(\ref{opp3}) & min & none --- {\it one\/} witness certifies \\
\end{tabular}
\end{center}
The same strategic move --- change the representation until its shape
fits the intended use --- carried each rung of the ladder.

\begin{thebibliography}{9}

\bibitem{ogryczak-tamir} W. Ogryczak and A. Tamir, Minimizing the sum of
the $k$ largest functions in linear time, {\it Information Processing
Letters} {\bf 85} (2003), 117--122.

\bibitem{boyd} S. Boyd and L. Vandenberghe, {\it Convex Optimization},
Cambridge University Press, 2004.

\bibitem{rockafellar-uryasev} R. T. Rockafellar and S. Uryasev,
Optimization of conditional value-at-risk, {\it Journal of Risk}
{\bf 2} (2000), 21--41.

\end{thebibliography}

\end{document}

